\centerline{\bf \Huge Регион ВсОШ}
\centerline{\bf \Huge 9 класс}

\begin{flushright}
        \scriptsize{}
\end{flushright}

\problem{1}
На окружности поставлены 143 точки, делящие её на равные дуги. Саша отметил одну из них. Первым ходом он отмечает соседнюю справа точку, вторым ходом он отмечает точку справа через одну от только что отмеченной, третьим - точку справа через две от только что отмеченной и т. д. На каком ходу впервые появится точка, отмеченная дважды?

\problem{2}
Можно ли на параболе $y=x^2$ отметить точки $A, B, C, D$, а на параболе $y=2 x^2$ точки $E, F, G, H$ так, чтобы четырёхугольники $A B C D$ и $E F G H$ оказались выпуклыми и равными?

\problem{3}
\textit{Близорукая} ладья бьёт все клетки своей строки и своего столбца, до которых можно дойти не более чем за 60 шагов, шагая каждый раз из клетки в соседнюю по стороне. Какое наибольшее число не бьющих друг друга близоруких ладей можно расставить на доске $100 \times 100$ ?

\problem{4}
Положительные числа $x, y, z$ таковы, что

$$
x y+y z+z x+2 x y z=1
$$


Докажите, что $4 x+y+z>2$.


\problem{5}
Точка $I_a$ - центр вневписанной окружности остроугольного треугольника $A B C$, касающейся стороны $B C$ в точке $X$, а точка $A^{\prime}$ диаметрально противоположна точке $A$ на описанной окружности этого треугольника. На отрезках $I_a X, B A^{\prime}, C A^{\prime}$ выбраны точки $Y, Z, T$ соответственно так, что $I_a Y=B Z=C T=r$, где $r$ - радиус вписанной окружности треугольника $A B C$. Докажите, что точки $X, Y, Z, T$ лежат на одной окружности.

